\begin{summarybox}
本章通过一个更广泛的动力学工具箱，替换了原先单一的“势阱”隐喻。
它的有据可依之处，不在于它已经用数学解决了行为问题，
而在于它给读者提供了一套有纪律的透镜：
用局部稳定性理解持续，用分离界理解路径敏感性，用极限环理解节律机制，用分岔理解阈值变化，用亚稳态理解灵活组织，用随机推理理解噪声与逃逸，用算子/控制视角理解受反馈支配的干预。
这是真实的数学严肃性、与生物学接触度以及设计精度的升级，
同时又不假装形式结构本身就已经构成对人类行动性的证明。
\end{summarybox}